\documentclass[11pt]{article}
\usepackage[a4paper,margin=28mm]{geometry}
\usepackage{amsmath,amssymb}
\usepackage{hyperref}
\title{Conjecture 00000008761 is false under its stated definition}
\author{Ziang Ni (\texttt{ziangni-sys})}
\date{10 October 2026}
\begin{document}
\maketitle
The English and Chinese statements define $SK_1$ as the kernel of
$\det:GL(R)\longrightarrow R^\times$ and assert its vanishing under a
nilradical condition. We disprove this vanishing clause using precisely
that definition.

Take the commutative field $R=\mathbb Q$. Its nilradical is zero: if
$x^m=0$ for $m>0$, the absence of zero divisors implies $x=0$.
Consequently it also satisfies any nilpotent-annihilation condition on
its nilradical: the zero ideal is annihilated already by the nilpotent
zero element. Passing to the reduced ring changes nothing.

Consider the actual invertible matrix
\[
 U=\begin{pmatrix}1&1\\0&1\end{pmatrix},\qquad
 U^{-1}=\begin{pmatrix}1&-1\\0&1\end{pmatrix}.
\]
Direct multiplication gives $UU^{-1}=U^{-1}U=I_2$, and
$\det U=1$. However $U\ne I_2$, because its $(1,2)$ entry is $1$
rather than $0$. Thus the determinant kernel is nontrivial.

This also works if $GL(R)$ denotes the stable general linear group.
For every $n\geq0$, the block diagonal matrix
$U_n=\operatorname{diag}(U,I_n)$ has determinant $1$ and retains the same
nonzero off-diagonal entry. Every standard stabilization is injective:
$\operatorname{diag}(A,1)=\operatorname{diag}(B,1)$ implies $A=B$
by restricting to the original block. Hence the class of $U$ never
becomes the identity in the direct limit. It is a nonidentity element of
the stable determinant kernel as well.

The claimed vanishing is therefore false. A false conjunct disproves
the whole conjecture; the subsequent claim about noncommutative
counterexamples need not be resolved.

\paragraph{Definition and formalization.}
The usual algebraic $K$-theory group $SK_1$ is instead a determinant
kernel \emph{after quotienting by elementary matrices}. Our matrix is
elementary and becomes trivial in that quotient. We make no claim
against a vanishing theorem for that different definition. The source
explicitly defines the unquotiented kernel, and that is the statement
refuted here.

The accompanying Lean project uses actual rational matrices and
the general linear group type. Its final theorem proves that
$\operatorname{nilradical}(\mathbb Q)=0$ and, for every $n$, there is a
nonidentity invertible $(2+n)$-dimensional matrix of determinant one.
The block indices $\mathrm{Fin}(2)\sqcup\mathrm{Fin}(n)$ express ordinary
stabilization. All axiom audits use only Lean's standard logical axioms.
\end{document}
